\documentclass{otengineering}

\typeout{OT7C-STARTER:OTENGINEERING}

\projectname{Starter Engineering Project}
\projectid{ENG-START}
\projectversion{0.1}
\projectstatus{\statusseed}
\projectcategory{Product Concept}
\projectstarted{\today}
\projectupdated{\today}

\begin{document}

\makeotengtitle

\section{Project Overview}

\begin{idea}
  Describe the product, mechanism, circuit, or system you want to explore.
\end{idea}

\begin{questionbox}
  What must be true for this idea to work in the real world?
\end{questionbox}

\section{First Decision}

\begin{decision}
  \entrydate{\today}
  \textbf{Decision:} Build a small paper or digital model before choosing final parts.

  \reason{A quick model exposes size, motion, and assembly problems early.}

  \tradeoff{The first model will not prove durability or manufacturing cost.}
\end{decision}

\section{First Experiment}

\begin{experiment}
  \entrydate{\today}
  \textbf{Objective:} Test the smallest risky assumption.

  \textbf{Method:} Record one measurement, one photo, or one clear observation.

  \successrating{5}
  \confidencerating{4}
  \pursuerating{Yes}
\end{experiment}

\section{Calculation}

\begin{calculation}
  \calcfield{Purpose}{Estimate whether the first design is physically plausible.}
  \calcfield{Inputs}{List dimensions, loads, voltages, currents, or materials.}
  \calcfield{Assumptions}{State what you are temporarily ignoring.}
  \calcfield{Calculation}{Write the first-order estimate here.}
  \calcfield{Interpretation}{Decide what the estimate says about the next prototype.}
\end{calculation}

\section{Sketch}

\sketchbox{Place a rough layout, mechanism sketch, or labelled block diagram here.}

\begin{futurememo}
  Keep the next step small enough to finish and test.
\end{futurememo}

\end{document}
